%==============================================================================
\chapter{Marco de referencia}
\label{cap:marco}
%==============================================================================

La construcción de un portafolio para el JSSP articula conceptos de la
programación de operaciones en planta, la programación por restricciones y la
selección automática de algoritmos mediante aprendizaje supervisado. A
continuación se revisan estos tres fundamentos y la revisión sistemática de
literatura que delimita la brecha abordada.

%------------------------------------------------------------------------------
\section{El Job Shop Scheduling Problem}
\label{sec:jssp}
%------------------------------------------------------------------------------
% Fuente: Anteproyecto, Cap. 6.1.1.

El Job Shop Scheduling Problem (JSSP) es uno de los problemas más estudiados
de la programación de tareas en entornos productivos. Consiste en asignar un
conjunto de $n$ trabajos, cada uno compuesto por una secuencia de
operaciones, a un conjunto de $m$ máquinas, de manera que cada operación se
procese en una máquina específica siguiendo un orden predefinido. El objetivo
clásico es minimizar el tiempo total de finalización (\textit{makespan})
\citep{pinedo2022}.

El problema es NP-duro, puesto que \citet{garey1976} demostraron que su versión
de decisión es NP-completa. Como cada máquina puede procesar sus $n$
operaciones en $n!$ órdenes distintos, el número de secuencias posibles crece
como $(n!)^m$, y por eso instancias de apenas $10 \times 10$, como
\texttt{ft10}, permanecieron sin resolver durante más de dos décadas
\citep{fisher1963, carlier1989}. En la práctica, los métodos exactos solo
garantizan el óptimo en instancias pequeñas o medianas dentro de un tiempo
razonable, y el desempeño de un mismo método varía mucho de una instancia a
otra, lo que motiva el uso de portafolios (Sección~\ref{sec:portafolios}).

Los métodos de solución suelen agruparse en tres familias
\citep{pinedo2022, strassl2022}. Los \textbf{métodos exactos} garantizan el óptimo si se
les da tiempo suficiente, e incluyen la ramificación y acotamiento específica
para el problema \citep{brucker1994}, la programación lineal entera mixta y la
programación por restricciones, que ha mostrado ser especialmente competitiva
en el JSSP \citep{ku2016}. Las \textbf{heurísticas constructivas} producen una
solución en un solo paso aplicando reglas de prioridad, como atender primero
la operación más corta \citep{panwalkar1977}. Entre ellas se encuentran el
algoritmo de \citet{giffler1960}, que genera programas activos, y el cuello de
botella móvil (\textit{shifting bottleneck}) de \citet{adams1988}. Son rápidas,
pero no garantizan la calidad de la solución. Las \textbf{metaheurísticas},
como la búsqueda tabú \citep{nowicki1996} o los algoritmos genéticos
\citep{yamada1992}, exploran el espacio de soluciones con estrategias
generales y suelen dar soluciones de muy buena calidad en instancias grandes,
aunque no prueban optimalidad. El portafolio de este trabajo se compone de
métodos exactos de programación por restricciones, y la heurística de
Giffler y Thompson se usa para acotar el problema
(Sección~\ref{sec:cotas}).

La investigación en el JSSP se apoya en instancias de referencia publicadas a
lo largo de varias décadas, entre las que destacan las de \citet{fisher1963},
\citet{lawrence1984}, \citet{adams1988}, \citet{applegate1991},
\citet{storer1992}, \citet{yamada1992}, \citet{taillard1993} y
\citet{demirkol1998}. Muchas se distribuyen a través de la OR-Library
\citep{beasley1990}, y sus cotas y mejores soluciones conocidas se actualizan
a medida que se publican nuevos resultados \citep{weise2019jssp}. Son las ocho
familias que se usan en este trabajo (Sección~\ref{sec:instancias}).

%------------------------------------------------------------------------------
\section{Programación por restricciones}
\label{sec:cp}
%------------------------------------------------------------------------------

En la programación por restricciones, el problema se describe de forma
declarativa mediante variables, dominios y restricciones, y un
\textit{solver} genérico se encarga de encontrar las soluciones
\citep{rossi2006}. Un problema de satisfacción de restricciones (CSP) se
define por una tripleta $(X, D, C)$ compuesta por variables $X$, un dominio
$D(x)$ de valores posibles para cada variable y restricciones $C$ que limitan
las combinaciones permitidas. Un problema de optimización con restricciones
(COP) añade una función objetivo. El JSSP con minimización del
\textit{makespan} es un COP cuyas variables son los tiempos de inicio de las
operaciones (Sección~\ref{sec:notacion}).

MiniZinc es un lenguaje de modelado de alto nivel, independiente del
\textit{solver} \citep{nethercote2007}, que separa el modelo (archivo
\texttt{.mzn}) de los datos de cada instancia (archivo \texttt{.dzn}). Antes
de resolver, el compilador traduce ambos a FlatZinc, el lenguaje de bajo nivel
que entienden los \textit{solvers}, y en esa traducción cada uno puede usar su
propia implementación de las restricciones globales, como
\texttt{disjunctive}, que impide que dos tareas usen a la vez el mismo
recurso. Gracias a esta separación, un único modelo del JSSP sirve para
ejecutar \textit{solvers} con técnicas muy distintas, y las diferencias de
desempeño se deben al \textit{solver}. MiniZinc es además el formato del
MiniZinc Challenge, la competencia anual de \textit{solvers} de restricciones
\citep{stuckey2014}.

Los \textit{solvers} clásicos combinan \textbf{propagación} y
\textbf{búsqueda} \citep{rossi2006}. La propagación elimina de los dominios
los valores que no pueden formar parte de ninguna solución. Por ejemplo, si
una operación no puede empezar antes del instante 30, su sucesora en el mismo
trabajo tampoco puede empezar antes de que la primera termine. Cuando la
propagación ya no avanza, la búsqueda elige una variable, le asigna un valor
o divide su dominio y vuelve a propagar, retrocediendo cuando encuentra una
contradicción. En optimización, cada solución encontrada impone que la
siguiente sea mejor, hasta probar que no existe una mejor. Gecode
\citep{gecode2019} es un \textit{solver} de este tipo, con propagadores
eficientes para muchas restricciones globales, pero, a diferencia de los
\textit{solvers} con generación perezosa de cláusulas, no aprende de los
fallos \citep{chu2011}.

La generación perezosa de cláusulas (\textit{Lazy Clause Generation}, LCG)
combina la propagación con el aprendizaje de cláusulas de los
\textit{solvers} SAT \citep{ohrimenko2009}. Cada reducción de un dominio se
registra con una explicación en forma de cláusula. Cuando se produce un fallo,
el \textit{solver} deriva de esas explicaciones una cláusula que resume la
causa del conflicto y la añade al problema, de modo que el mismo fallo no se
repite en otras ramas. Chuffed \citep{chu2011} implementa esta idea, que es
especialmente eficaz en problemas de \textit{scheduling}
\citep{schutt2011}.

CP-SAT, el \textit{solver} de restricciones de OR-Tools \citep{perron2025},
combina en un mismo motor un núcleo SAT con generación perezosa de cláusulas,
relajaciones de programación lineal que aportan cotas al objetivo, técnicas
de presolución y búsqueda local, y por defecto ejecuta en paralelo varias
estrategias de búsqueda \citep{perron2023}, aunque en este trabajo se usó con un único hilo, en
igualdad de condiciones con los demás \textit{solvers}
(Sección~\ref{sec:ejecucion_solvers}). Por sus resultados en el MiniZinc
Challenge y en \textit{scheduling}, se usa aquí como referencia del estado
del arte.

%------------------------------------------------------------------------------
\section{El problema de selección de algoritmos}
\label{sec:asp}
%------------------------------------------------------------------------------
% Fuente: Anteproyecto, Cap. 6.1.3.

\citet{rice1976} planteó el problema de selección de algoritmos como el de
determinar qué algoritmo, entre un conjunto disponible, es el más adecuado
para una instancia específica. Su modelo distingue cuatro espacios: el de
problemas $\mathcal{P}$, que contiene las instancias; el de características
$\mathcal{F}$, donde cada instancia $i$ se representa por un vector $f(i)$; el
de algoritmos $\mathcal{A}$, y el de desempeño $\mathcal{Y}$, donde $y(a, i)$
mide qué tan bien resuelve $a$ la instancia $i$. Se trata de encontrar una
función $S: \mathcal{F} \rightarrow \mathcal{A}$ que elija el algoritmo con
mejor desempeño esperado, y en la práctica $S$ se aprende de instancias cuyo
desempeño con cada algoritmo ya se conoce \citep{kotthoff2014,
kerschke2019}. Iniciativas como COSEAL promueven metodologías y conjuntos de
datos comunes para comparar estrategias de selección \citep{coseal2024}, y
las competencias de selección de algoritmos han fijado métricas comunes
\citep{lindauer2019}. Este trabajo extiende ese planteamiento para que la
función produzca un \textit{schedule} de ejecución completo
(Sección~\ref{sec:formulacion_schedule}).

La calidad de un selector depende de que las características describan bien
lo que hace difícil una instancia. Las \textbf{características genéricas} se
calculan sobre la representación que recibe el \textit{solver}, sin conocer
el dominio. En programación por restricciones, \texttt{mzn2feat}
\citep{amadini2014mzn2feat} extrae características del modelo FlatZinc (95 en
su versión 1.0, \citealp{amadini2015sunnycp}), y CPHydra
\citep{omahony2008cphydra} usa un conjunto más reducido de características
del mismo tipo. Sirven para
cualquier problema, pero dependen de cómo esté escrito el modelo. Las
\textbf{características específicas} usan conocimiento del dominio, como
ocurre en el JSSP con la relación entre trabajos y máquinas, la variabilidad
de los tiempos de proceso, la carga de las máquinas y la estructura de las
rutas \citep{mirshekarian2016, strassl2022}. Este trabajo combina ambos tipos
(Sección~\ref{sec:features}).

%------------------------------------------------------------------------------
\section{Portafolios de algoritmos}
\label{sec:portafolios}
%------------------------------------------------------------------------------
% Fuente: Anteproyecto, Cap. 6.1.2.

Un portafolio de algoritmos es un conjunto de métodos para un mismo tipo de
problema, cada uno con ventajas y debilidades propias, del que se elige o
combina el más apropiado según la instancia \citep{gomes2001, kotthoff2014}. Su
motivación es que ningún método es el mejor en todos los casos, como lo
expresa el teorema \textit{No Free Lunch} \citep{wolpert1997}.

La estrategia más común, que aquí se llama \textit{Winner Takes All} (WTA)
como en el anteproyecto, elige un único algoritmo por instancia y le dedica
todo el tiempo disponible. El selector puede ser un clasificador que predice
el mejor algoritmo o un conjunto de modelos de regresión que predicen el
desempeño de cada uno \citep{xu2008satzilla, kotthoff2014}. Si el algoritmo
elegido falla, ya no queda tiempo para otro, y ese riesgo crece cuando las
predicciones son poco fiables.

Un \textit{schedule} secuencial de ejecución reparte el presupuesto $T$ entre
varios algoritmos que se ejecutan uno tras otro. A cambio de dar menos tiempo
a cada uno, reduce el riesgo de una mala selección, y suele bastar porque
muchas instancias se resuelven rápido con el algoritmo adecuado
\citep{omahony2008cphydra, amadini2014sunny}. Construirlo exige decidir qué
algoritmos incluir, cuánto tiempo darle a cada uno y en qué orden ejecutarlos,
que son las tres decisiones de la pregunta de investigación de este trabajo.

Entre los sistemas representativos, SATzilla \citep{xu2008satzilla} es un
portafolio WTA para SAT que predice el tiempo de cada \textit{solver} y elige
el menor, tras ejecutar brevemente unos \textit{presolvers}. Hydra
\citep{xu2010hydra} construye el portafolio configurando versiones de un
mismo algoritmo parametrizado que se complementen entre sí. CPHydra
\citep{omahony2008cphydra}, el primer portafolio de \textit{solvers} para
programación por restricciones \citep{amadini2015sunnycp}, busca las instancias más parecidas a la nueva
y calcula el \textit{schedule} que habría resuelto más de ellas. aspeed
\citep{hoos2015aspeed} formula la construcción del \textit{schedule} como un
problema de optimización resuelto con programación de conjuntos de
respuestas.

SUNNY \citep{amadini2014sunny}, la referencia metodológica de este trabajo,
simplifica la idea de CPHydra sin fase de entrenamiento. Para una instancia
nueva busca sus $k$ vecinos más cercanos en una base de conocimiento y
selecciona el menor subconjunto de \textit{solvers} que resuelve el máximo
número de ellos. Divide el tiempo en franjas iguales y asigna a cada
\textit{solver} tantas como vecinos resuelve. Las franjas de los vecinos que
ninguno resuelve pasan a un \textit{solver} de respaldo, y el orden de
ejecución sigue el tiempo medio en el vecindario. Para optimización, SUNNY-CP
\citep{amadini2015sunnycp} reemplaza el conteo de instancias resueltas por una
puntuación que premia probar el óptimo y la calidad de la solución.

%------------------------------------------------------------------------------
\section{Aprendizaje automático para predecir el desempeño de algoritmos}
\label{sec:ml}
%------------------------------------------------------------------------------

Los modelos que predicen el desempeño de un algoritmo pueden ser de
regresión, si estiman una magnitud continua como el tiempo de ejecución, o de
clasificación, si predicen una categoría como ``resuelve antes del límite''.
Predecir el tiempo a partir de características es una técnica consolidada
\citep{hutter2014}. Entre los modelos más usados están los conjuntos de
árboles, tales como el \textbf{bosque aleatorio} \citep{breiman2001}, que
promedia muchos árboles entrenados sobre muestras aleatorias y es uno de los
más precisos para esta tarea \citep{hutter2014}, y el \textit{gradient
boosting} \citep{friedman2001}, que construye los árboles en secuencia, cada
uno ajustado a los errores de los anteriores. Los tiempos suelen estar
\textbf{censurados}, porque cuando un algoritmo agota el límite solo se sabe
que su tiempo real es mayor, por lo cual conviene predecir su logaritmo y
tratar con cuidado esos valores censurados \citep{hutter2014}.

El método de $k$ vecinos más cercanos ($k$-NN) \citep{cover1967} predice el
comportamiento de una instancia a partir de las $k$ más parecidas según una
distancia en el espacio de características. Como su ``modelo'' es la propia
base de instancias, es sensible a la escala de las características, que se
normalizan antes de calcular distancias, y a las características
irrelevantes o redundantes. SUNNY usa $k$-NN para construir el vecindario, y
en sus experimentos valores de $k$ entre 5 y 20 dan resultados similares
\citep{amadini2014sunny}.

\subsection{Métricas para evaluar portafolios}
\label{sec:metricas_portafolios}

Un portafolio se evalúa frente a dos referencias \citep{kerschke2019}. El
\textbf{SBS} (\textit{Single Best Solver}) es el algoritmo con mejor
desempeño promedio, lo que se obtendría usando siempre el mismo, de modo que un
portafolio solo aporta si logra superarlo. El \textbf{VBS} (\textit{Virtual
Best Solver}) es un oráculo que elige en cada instancia el mejor algoritmo, y
marca el límite de lo que puede lograr cualquier portafolio con esos
algoritmos. La diferencia entre ambos mide el margen disponible, y una medida
habitual es la fracción de ese margen que el portafolio logra cerrar
\citep{kerschke2019, lindauer2019}.

Las métricas de desempeño más usadas son el número de instancias resueltas
(en optimización, con el óptimo probado), el \textbf{PAR10}, que promedia el
tiempo contando cada instancia no resuelta como $10 \cdot T$
\citep{lindauer2019}, el tiempo hasta probar el óptimo y el
\textbf{\textit{gap}} frente a la mejor solución conocida
(Sección~\ref{sec:bks_gap}), que mide la calidad cuando el óptimo no se
prueba. SUNNY-CP usa además una puntuación que vale 1 si se prueba el óptimo,
entre 0.25 y 0.75 según la calidad de la mejor solución y 0 si no se
encuentra ninguna \citep{amadini2015sunnycp}.

%------------------------------------------------------------------------------
\section{Estado del arte}
\label{sec:estado_arte}
%------------------------------------------------------------------------------

El estado del arte se construyó mediante una revisión sistemática de
literatura que siguió la guía de \citet{carrera2022}. El protocolo PICOC
(población, intervención, comparación, resultado y contexto) delimitó el
alcance de la búsqueda (Tabla~\ref{tab:picoc}).

\begin{table}[H]
  \centering
  \small
  \singlespacing
  \caption{Protocolo PICOC de la revisión sistemática.}
  \label{tab:picoc}
  \begin{tabularx}{\textwidth}{>{\raggedright\arraybackslash}p{0.16\textwidth}
                               >{\raggedright\arraybackslash}p{0.22\textwidth}
                               >{\raggedright\arraybackslash}X}
    \toprule
    \textbf{Criterio} & \textbf{Término principal} & \textbf{Términos relacionados} \\
    \midrule
    Población    & Job Shop Scheduling Problem & JSSP, \textit{makespan},
      \textit{scheduling} NP-duro. Se restringe a la versión clásica. \\
    Intervención & Portafolio de algoritmos & \textit{Algorithm portfolio},
      \textit{schedule} de ejecución, \textit{Winner Takes All}. \\
    Comparación  & Algoritmo único & Búsqueda tabú, algoritmos genéticos,
      recocido simulado, programación lineal entera mixta. \\
    Resultado    & Selección automática & Problema de selección de
      algoritmos, características de instancias, selección por instancia. \\
    Contexto     & \textit{Machine Learning} & Aprendizaje supervisado,
      bosques aleatorios, máquinas de soporte vectorial, redes neuronales. \\
    \bottomrule
  \end{tabularx}
\end{table}

Del protocolo se derivaron tres preguntas de investigación, que guiaron la
extracción de información de cada artículo:

\begin{description}
  \item[RQ1.] ¿Cuáles son los enfoques de portafolios de algoritmos más
    efectivos aplicados al JSSP, comparados con la estrategia
    \textit{Winner Takes All}?
  \item[RQ2.] ¿Qué características de las instancias del JSSP son más
    relevantes para predecir el desempeño de los algoritmos de resolución?
  \item[RQ3.] ¿Cómo se aplican las técnicas de \textit{Machine Learning} para
    construir \textit{schedules} de ejecución de portafolios en problemas de
    optimización combinatoria NP-duros?
\end{description}

La búsqueda priorizó bases de datos y repositorios de acceso libre, entre
ellos arXiv y dblp, y se limitó a artículos en inglés publicados entre 2001 y 2025 en revistas o conferencias
con revisión por pares y relacionados con al menos dos criterios PICOC. Se
excluyeron los libros sin acceso digital, los trabajos en otros idiomas que no
fueran inglés o español y los centrados exclusivamente en el Flexible Job Shop
o en variantes multiobjetivo. La revisión seleccionó 12 trabajos, organizados
en tres bloques temáticos (Tabla~\ref{tab:slr}). Durante el desarrollo se
incorporaron otras referencias, identificadas a partir de estos trabajos y de
las necesidades de cada fase, sin aplicar de nuevo el protocolo.

\begin{table}[H]
  \centering
  \small
  \singlespacing
  \caption{Trabajos seleccionados en la revisión sistemática.}
  \label{tab:slr}
  \begin{tabularx}{\textwidth}{>{\raggedright\arraybackslash}p{0.30\textwidth}
                               >{\raggedright\arraybackslash}p{0.17\textwidth}
                               >{\raggedright\arraybackslash}X}
    \toprule
    \textbf{Trabajo} & \textbf{Bloque} & \textbf{Aporte al proyecto} \\
    \midrule
    \citet{pinedo2022} & JSSP & Marco formal del JSSP, \textit{makespan} y
      métodos de solución. \\
    \citet{reijnen2023} & JSSP & Plataforma unificada de instancias y métodos
      de referencia para el JSSP y sus variantes. \\
    \citet{xu2025learn} & JSSP / ML & Compara programación genética y
      aprendizaje por refuerzo para aprender heurísticas del JSSP. \\
    \citet{smit2024gnn} & JSSP / ML & Revisión de redes neuronales de grafos
      aplicadas al JSSP. \\
    \citet{gomes2001} & Portafolios & Trabajo fundacional sobre portafolios de
      algoritmos y la variabilidad de su desempeño. \\
    \citet{xu2008satzilla} & Portafolios & Portafolio de selección por
      instancia; arquitectura de referencia. \\
    \citet{streeter2008} & Portafolios & Diseño de portafolios y
      \textit{schedules} de ejecución con reinicio o suspensión. \\
    \citet{amadini2014sunny} & Portafolios & \textit{Schedule} de ejecución
      basado en vecinos más cercanos; referencia metodológica del trabajo. \\
    \citet{omahony2008cphydra} & Portafolios & Portafolio con
      \textit{schedule} de ejecución para programación por restricciones,
      basado en las instancias más parecidas; el primero de su tipo según
      \citet{amadini2015sunnycp}. \\
    \citet{kotthoff2014} & Selección & Clasificación de las técnicas de
      aprendizaje para la selección de algoritmos. \\
    \citet{kerschke2019} & Selección & Revisión de la selección automática de
      algoritmos y de la caracterización de instancias. \\
    \citet{bengio2021} & ML en optimización & Marco metodológico para
      integrar aprendizaje automático y optimización combinatoria. \\
    \bottomrule
  \end{tabularx}
\end{table}
% Nota: el informe de la revisión anotó como "Cunha et al. (2025)", "Hao et al.
% (2024)" y "Xu et al. (2012)" los trabajos de Xu et al. (2025), Smit et al.
% (2024) y Streeter y Smith (2008); se citan con sus autores reales. El informe
% cita Kotthoff (2016), versión en libro del survey; aquí se usa Kotthoff (2014).

\subsection{Síntesis de la revisión}
\label{sec:evolucion_jssp}
\label{sec:ml_co}
% Fuente: Anteproyecto, Cap. 6.2.

Los métodos para el JSSP abarcan desde los primeros algoritmos heurísticos
\citep{giffler1960} y los métodos exactos basados en ramificación y
acotamiento y programación entera hasta metaheurísticas y métodos híbridos
que incorporan aprendizaje \citep{pinedo2022}. En paralelo, los portafolios de
algoritmos se consolidaron en satisfacción de restricciones, planificación y
otros problemas NP-duros, con sistemas como SATzilla \citep{xu2008satzilla}
e Hydra \citep{xu2010hydra}, que caracterizan las instancias y predicen el
algoritmo o la configuración más apropiada. \citet{kotthoff2014} describe
portafolios de algoritmos y de configuraciones, y señala la ejecución paralela
como una línea de investigación en curso. En optimización combinatoria, el aprendizaje automático se ha usado para
predecir el rendimiento de los algoritmos, guiar búsquedas heurísticas y
ajustar parámetros de optimización \citep{bengio2021}, y en el JSSP en
particular para aprender heurísticas y reglas de despacho
\citep{smit2024gnn, xu2025learn}.

%------------------------------------------------------------------------------
\section{Síntesis y brecha identificada}
\label{sec:brecha}
%------------------------------------------------------------------------------

De la literatura revisada se desprende que el desempeño de los métodos para el
JSSP varía de manera pronunciada entre instancias. Los portafolios aprovechan
esa variabilidad, aunque la mayoría elige un único algoritmo por instancia y
queda expuesta a errores de selección que los \textit{schedules} secuenciales
como CPHydra, SUNNY o aspeed mitigan repartiendo el tiempo. Asimismo, el
aprendizaje automático, en particular los vecinos más cercanos y los conjuntos
de árboles, es la herramienta habitual para conectar las características de una
instancia con el desempeño esperado de cada algoritmo.

En los trabajos revisados no se encontró un portafolio por instancia ni un
\textit{schedule} de ejecución al estilo de SATzilla o SUNNY aplicado al JSSP
clásico, ya que los portafolios secuenciales de programación por restricciones
se han evaluado sobre colecciones heterogéneas con características genéricas.
Las revisiones sobre aprendizaje automático en el JSSP se concentran en la
generación de heurísticas y políticas de despacho, con escasa atención a la
selección de algoritmos, y tampoco se encontró un conjunto de datos público con
el desempeño de varios métodos sobre el JSSP, análogo a ASlib
\citep{bischl2016aslib}.

Este trabajo aborda esas brechas aplicando y comparando la selección por
instancia y los \textit{schedules} de ejecución en el JSSP clásico, con
\textit{solvers} de restricciones que comparten un modelo en MiniZinc, y
construye una base de conocimiento con el desempeño de cada configuración en
las familias de instancias de referencia del problema
(Capítulo~\ref{cap:caracterizacion}). Con ella estudia si un
\textit{schedule} aporta frente a la selección de una sola configuración
cuando el portafolio se especializa en un solo problema, y qué tan
informativas son las características genéricas y las propias del JSSP.
